\section{Lecture 20 (November 25th)}
\begin{rmk}
Let $f$ be a entire function. The zero set of $f$ is deonted as $Z_{f}$. If $A$ is a infinite subset of ${\bm C}$ and $Z_{f}=A$ for some entire function $f$, then $A$ must not have a limit point. If $A=\{a_{n}\}\subset {\bm C}$ with $\lim _{n\rightarrow \infty }|a_{n}|=\infty $, then is there an entire function $f$ satisfying $A=Z_{f}$? The answer is yes.
\end{rmk}
\vspace{2ex}
\begin{ex}
Let $A=\{a_1,\ldots ,a_{n}\}$ be a finite zero set. Then $f(z)=(z-a_1)\ldots (z-a_{n})$ and $A=Z_{f}$. 
\end{ex}
\vspace{2ex}
\begin{prop}
If there are two entire functions $f$ and $g$ such that $Z_{f}=Z_{g}$, then $\dfrac{f}{g}$ is entire with no zeros. Since ${\bm C}$ is simply connected, there exists an entire function $h$ such that $\dfrac{f}{g}=e^{h}$. That is, $f=ge^{h}$. 
\end{prop}
\vspace{2ex}
\begin{ex}
An entire function $f$ such that $Z_{f}={\bm Z}$.
\[f(z)=\sin \pi z\]
How do we find an entire function $f$ such that $Z_{f}=\{2,3,4,5,6,\ldots \}$? Let's use the notation
\[\prod^{N}_{n=1}a_{n}=a_1a_2\ldots a_{N}\]
and write
\[f(z)=\Big(1-\dfrac{z}{2}\Big)\Big(1-\dfrac{z}{3}\Big)\Big(1-\dfrac{3}{4}\Big)\]
This function is zero at $1$, and our conjecture is false. The idea is to find a multiplier $e^{g(z)}$ to take this property away!
\[f(z)=\prod^{\infty }_{n=2}\Big(1-\dfrac{z}{n}\Big)e^{g(z)}\]
\end{ex}
\vspace{2ex}
\begin{lem}
If $a_{n}\geq 0$, then 
\[\prod^{\infty }_{n=1}(1+a_{n})\]
converges if and only if 
\[\sum ^{\infty }_{n=1}a_{n}\]
converges. 
\end{lem}
\vspace{2ex}
\begin{proof}
If $x\geq 0$,
\[e^{x}=1+x+\dfrac{x^2}{2!}+\dfrac{x^3}{3!}\geq 1+x\]
\end{proof}
\vspace{2ex}
\begin{proof}
\[\sum ^{N}_{n=1}a_{n}\leq \prod^{N}_{n=1}(1+a_{n})\leq \prod^{N}_{n=1}e^{a_{n}}=e^{\sum ^{N}_{n=1}a_{n}}\]
using the comparison test. 
\end{proof}
\vspace{2ex}
\begin{thm}
(Big theorem I) Let $a_{n}\in {\bm C}$. If $\sum ^{\infty }_{n=1}|a_{n}|<\infty $, then $\prod^{\infty }_{n=1}(1+a_{n})$ converges and if $\prod_{n=1}^{\infty }(1+a_{n})=0$ then $a_{n}=-1$ for some $n\in {\bm N}$. 
\end{thm}
\vspace{2ex}
\begin{proof}
Since $\sum ^{\infty }_{n=1}|a_{n}|$ converges, there is $K\in {\bm N}$ such that if $n\geq K$, then $|a_{n}|<\dfrac{1}{2}$. Then 
\[\prod^{N}_{n=k} (1+a_{n})=\prod^{N}_{n=k}e^{\log (1+a_{n})}=e^{\sum ^{N}_{n=k}\log (1+a_{n})}\]
Notice that
\[\sum ^{N}_{n=k}|\log (1+a_{n})|\leq \sum ^{N}_{n=k}2|a_{n}|\]
by lemma. Thus,
\[\sum _{n=k}^{\infty }\log(1+a_{n})\]
converges to some $A\in {\bm C}$. 
\[\prod^{N}_{n=k}(1+a_{n})=e^{A}\ne 0\]
\end{proof}
\vspace{2ex}
\begin{lem}
We denote $\log z$ for $z\ne 0$ and $-\pi <\mathrm{arg}\,z<\pi $. 
\[\log z=\ln |z|+i\mathrm{arg}\,z\]
If $|z|<\dfrac{1}{2}$, then $\log(1+z)$ satisfies $|\log(1+z)|\leq 2|z|$.
\end{lem}
\vspace{2ex}
\begin{proof}
Since $\log (1+z)$ is analytic in $|z|<1$, 
\[\log (1+z)=z-\dfrac{z^2}{2}+\dfrac{z^3}{3}-\dfrac{z^{4}}{4}\]
for $|z|<1$. Thus, if $|z|<\dfrac{1}{2}$, \[|\log (1+z)|=\Big|z-\dfrac{z^2}{2}+\dfrac{z^3}{3}+\ldots \Big|\leq |z|\Big(1+\dfrac{1}{2}+\dfrac{1}{4}+\ldots \Big)=2|z|\]
\end{proof}
\vspace{2ex}
\begin{thm}
(Big theorem II) Let $\{f_{n}\}$ be analytic in a domain $D$. If there is $\{M_{n}\}$ with $M_{n}\geq 0$ satisfying $\sum ^{\infty }_{n=1}M_{n}<\infty $ and $|f_{n}(z)-1|\leq M_{n}$ for all $z\in D$, then
\[f(z)=\prod^{\infty }_{n=1}f_{n}(z)\]
is analytic in $D$ since $\prod_{n=1}^{N}f_{n}(z)$ converges uniformly on every compact subset. 
\end{thm}
\vspace{2ex}
\begin{proof}
There exists $n\geq k$ such that $|f_{n}(z)-1|<\dfrac{1}{2}$ for all $z\in D$.
\[\prod^{N}_{n=k}f_{n}(z)=\prod_{n=k}^{N}e^{\log (f_{n}(z))}=e^{\sum \log(f_{n}(z))}\leq e^{\sum M_{n}}\]
as $|\log (f_{n}(z)-1)|\leq M_{n}$. 
\end{proof}
\vspace{2ex}
\begin{lem}
If $f_{n}$ converges uniformly and $g$ is uniformly continuous,
\[g\circ f_{n}\]
converges uniformly.
\end{lem}
\vspace{2ex}
\begin{ex}
Prove that $\prod^{\infty }_{n=1}\Big(1-\dfrac{z}{n^2}\Big)$ is an entire function with $Z_{f}=\{n^2\;|\;n\in {\bm N}\}$. 
\end{ex}
\vspace{2ex}
\begin{proof}
Let $D=D(0,R)$. Take $f_{n}(z)=z/n^2$ and $M_{n}=R/n^2$ we can conclude that $f(z)$ is analytic in $D$. 
\end{proof}
\vspace{2ex}

